% ==========================================================
%  第 1 章 drill：结构雷达闪电战 + 五问连环母题
% ------------------------------------------------------------
%  纯模块化片段：不含 documentclass / document；
%  编译入口：_check_ch01.tex（preamble + versions + 本文件）。
%  只使用 versions.tex 与 common/ 已提供的环境（warmupbox、
%  advancedonly、\ifteacher、\tk、\logotag 等），不新增任何
%  \newcommand / \newenvironment。
%  数值核验（sympy 实算，见工作记录）：
%    母题一 f=x^3+x^2+x+1 的偶部 x^2+1、奇部 x^3+x，恒等式与
%    特殊值 f(1)=4,f(-1)=0,g(±1)=2,h(±1)=±2 全部成立；
%    母题二 x1=nsolve(e^x+x-4)≈1.0737289，x2=3-x1≈1.9262711，
%    残差 ln(x2+1)+1+x2-4=0，x1+x2=3、y1+y2=5；
%    P'(e-1,2) 代入 C2 精确为 0；xy=1 配 y=x+c 的韦达根和 =-c。
% ==========================================================

% ---------------------------------------------------------
%  第一部分：五分钟结构雷达闪电战
% ---------------------------------------------------------

\begin{warmupbox}{五分钟结构雷达闪电战：只练核心一步变形}
规则：每题\textbf{只写关键一步}（不是解完整题）。先读「雷达触发」，在脑子里把这一步放回整条链里，再动笔；答案处以下划线留白（学生版显示等宽下划线，教师版显红）。

\begin{enumerate}[label=(\arabic*), itemsep=5pt]
\item \textbf{奇偶正交分解}：把 $f(x)=x^3+2x^2+x+1$ 分解为「偶部＋奇部」。\\
{\small\bfseries\textcolor{third}{雷达触发：}}先写 $f(-x)$：偶次项不变、奇次项变号；相加除 $2$ 留偶部，相减除 $2$ 留奇部。\\
答案：偶部 \tk{2x^2+1}，奇部 \tk{x^3+x}。

\item \textbf{关于 $y$ 轴的镜像反射}：写出曲线 $y=x^3-x$ 关于 $y$ 轴对称后的表达式。\\
{\small\bfseries\textcolor{third}{雷达触发：}}哪个坐标变号，就把哪个变量取反；横坐标变 $\Rightarrow x\to-x$，只替换自变量。\\
答案：\tk{y=-x^3+x}（即 $f(-x)$）。

\item \textbf{$180^\circ$ 中心旋转}：点 $(a,b)$ 绕原点旋转 $180^\circ$ 的像是 \tk{(-a,-b)}；曲线 $y=x^2+1$ 旋转后的表达式是 \tk{y=-x^2-1}。\\
{\small\bfseries\textcolor{third}{雷达触发：}}两次镜像（先 $y$ 轴、再 $x$ 轴）合成中心旋转——两个坐标都取反，再解出新 $y$。

\item \textbf{刚体平移（顶点追踪）}：抛物线 $y=x^2-2x+3=(x-1)^2+2$ 的顶点移到 $(4,3)$ 后，新抛物线是 \tk{y=(x-4)^2+3}。\\
{\small\bfseries\textcolor{third}{雷达触发：}}平移不改开口与形状，只追顶点：$(1,2)\to(4,3)$，即右移 $3$、上移 $1$。

\item \textbf{双对称推周期}：$f$ 的图象关于直线 $x=1$ 与 $x=4$ 都对称，则最小正周期 $T=\tk{6}$。\\
{\small\bfseries\textcolor{third}{雷达触发：}}连续两次平行镜像反射合成一次纯平移，平移量 $2|b-a|$ 就是周期（见本章方法块）。

\item \textbf{单值性垂线检验}：点集 $\{(x,y)\mid x^2+y^2=1\}$ 是不是某个函数的图象？\\
{\small\bfseries\textcolor{third}{雷达触发：}}证伪只需找一条截出两个点的竖直线——$x=0$ 交于 $(0,1),(0,-1)$。\\
答案：\tk{\text{不是}}（竖直线 $x=0$ 与它有两个交点）。
\end{enumerate}
\end{warmupbox}

% ---------------------------------------------------------
%  第二部分：解答即笔记 · 五问连环母题
% ---------------------------------------------------------

\begin{tcolorbox}[
  breakable, enhanced, sharp corners,
  colback=structurecolor!3, colframe=structurecolor!70!black, boxrule=0.6pt,
  title={\bfseries 解答即笔记·五问连环母题一：对称即不变性——三种变换下的「什么不变、什么改变」},
  coltitle=black, fonttitle=\small,
  attach boxed title to top left={xshift=2mm, yshift=-2mm},
  boxed title style={sharp corners, colback=white,
    colframe=structurecolor!70!black, boxrule=0.4pt},
  top=7pt]

试验品 $f(x)=x^3+x^2+x+1$（定义域 $\mathbb{R}$）。三个变换：镜像反射 $M_y:(x,y)\mapsto(-x,y)$（表达式 $f(x)\to f(-x)$）；$180^\circ$ 中心旋转 $R_O:(x,y)\mapsto(-x,-y)$（表达式 $f(x)\to-f(-x)$）；刚体平移 $T:(x,y)\mapsto(x+h,y+k)$（表达式 $y-k=f(x-h)$）。

\begin{enumerate}[label=(\arabic*), itemsep=4pt]
\item \textbf{具象感知与画图}：先算 $f(1)=\tk{4}$，$f(-1)=\tk{0}$；再把 $f$ 与其偶部、奇部画在同一坐标系里，观察三条曲线各自满足哪种对称。

\begin{center}
\begin{tikzpicture}[>=Stealth, semithick, font=\footnotesize, x=1.6cm, y=0.9cm]
  \draw[gray!22] (-1.7,-3) grid (1.7,5.4);
  \draw[->, gray] (-1.9,0)--(1.95,0) node[below right] {$x$};
  \draw[->, gray] (0,-3.2)--(0,5.7) node[above left] {$y$};
  \foreach \x in {-1,1}{ \draw (\x,0.08)--(\x,-0.08) node[below=1pt] {\scriptsize \x}; }
  \foreach \y in {-2,2,4}{ \draw (0.08,\y)--(-0.08,\y) node[left=1pt] {\scriptsize \y}; }
  \draw[main!70!black, very thick, domain=-1.25:1.25, samples=60, smooth] plot(\x,{(\x)*(\x)+1});
  \draw[second, very thick, domain=-1.25:1.25, samples=60, smooth] plot(\x,{(\x)*(\x)*(\x)+\x});
  \draw[third, very thick, dashed, domain=-1.25:1.25, samples=60, smooth] plot(\x,{(\x)*(\x)*(\x)+(\x)*(\x)+\x+1});
  \node[main!80!black] at (1.62,2.55) {$g=x^2+1$};
  \node[second!80!black] at (1.62,-1.95) {$h=x^3+x$};
  \node[third!70!black] at (-1.62,4.62) {$f=g+h$};
  \ifteacher
    \fill[main] (1,2) circle (1.5pt);
    \fill[main] (-1,2) circle (1.5pt);
    \fill[second] (1,2) circle (1.5pt);
    \fill[second] (-1,-2) circle (1.5pt);
    \fill[third] (1,4) circle (1.5pt) node[left] {$(1,4)$};
    \fill[third] (-1,0) circle (1.5pt) node[right] {$(-1,0)$};
    \node[main!80!black, align=left, font=\scriptsize] at (0,5.25) {偶部：关于 $y$ 轴镜像重合};
    \node[second!80!black, align=left, font=\scriptsize] at (0,-2.9) {奇部：绕原点转 $180^\circ$ 重合};
  \else
    \node[third] at (0,-2.9) {描点作图：哪条关于 $y$ 轴对称？哪条绕原点旋转 $180^\circ$ 重合？$f$ 呢？};
  \fi
\end{tikzpicture}
\end{center}

\ifteacher
{\small\textcolor{second}{红笔：$x=1$ 处 $g(1)+h(1)=2+2=4=f(1)$；$x=-1$ 处 $2+(-2)=0=f(-1)$。$g$ 在 $x$ 与 $-x$ 处等高（镜像不变）；$h$ 在 $x$ 与 $-x$ 处等值反号（旋转不变：$(x,h(x))$ 与 $(-x,-h(x))$ 同在曲线上）。$f$ 两种对称都不满足——但分解之后，各部分都找到了自己的对称性。}}
\else
\noteblank[3ex]
\fi
\logotag{具象感知}

\item \textbf{考纲满分通法}：按定义证明「任意函数＝偶部＋奇部」，并套用到试验品。前提：$f$ 的定义域 $D$ 关于原点对称（$\forall x\in D,\ -x\in D$）——这是「不变性」的舞台，缺了它连 $f(-x)$ 都并非处处有意义。

设
\[
g(x)=\tk{\frac{f(x)+f(-x)}{2}},\qquad h(x)=\tk{\frac{f(x)-f(-x)}{2}}.
\]
则
\[
g(-x)=\tk{\frac{f(-x)+f(x)}{2}=g(x)}\ (\text{偶}),\qquad
h(-x)=\tk{\frac{f(-x)-f(x)}{2}=-h(x)}\ (\text{奇}),
\]
且 $g+h=f$。套用：$f(-x)=\tk{-x^3+x^2-x+1}$，故 $g=\tk{x^2+1}$，$h=\tk{x^3+x}$。
\logotag{考纲通法}

\item \textbf{高观点降维秒杀}：「对称即不变性」升华卡——变换 $T$ 作用在图象上，$T(C)=C$ 的部分就是不变量。查表即可秒判，不必逐条重算。

\begin{center}\small\renewcommand{\arraystretch}{1.35}
\begin{tabular}{p{2.7cm}|p{2.5cm}|p{4.5cm}|p{4.5cm}}
\toprule
变换（点集规则） & 表达式 & 什么不变 & 什么改变 \\
\midrule
镜像 $M_y$：$(x,y)\mapsto(-x,y)$ & $f(x)\to f(-x)$ & 值域；$|f|$ 随 $|x|$ 的分布；奇偶身份；周期 & 零点 $x_0\to-x_0$；对称轴 $a\to-a$；对称区间上的单调方向 \\
旋转 $R_O$：$(x,y)\mapsto(-x,-y)$ & $f(x)\to-f(-x)$ & $|f|$ 分布；奇偶身份；周期；单调方向；最值的绝对值 & 值域上下翻转（最值 $\max\leftrightarrow\min$）；零点与对称中心 $(a,b)\to(-a,-b)$ \\
平移 $T_{(h,k)}$：$(x,y)\mapsto(x+h,y+k)$ & $y-k=f(x-h)$ & 形状与单调方向；周期；极值大小；值域宽度 & 对称轴／中心的位置；奇偶身份（一般被破坏）；零点位置 \\
\bottomrule
\end{tabular}
\end{center}

\textbf{正交分解}：任意 $f=\dfrac{f(x)+f(-x)}{2}+\dfrac{f(x)-f(-x)}{2}$——偶部是 $M_y$ 的不变量，奇部是 $R_O$ 的不变量；偶×奇＝奇，两部分叠加时各保各的对称性、互不干扰（「正交」）。回代验算：$x=1$ 处 $2+2=4=f(1)$，$x=-1$ 处 $2-2=0=f(-1)$。

\textbf{挂上本章}：图象变换图谱的一级（平移／轴对称／中心对称）就是这张表的前两列；方法块「双对称推周期」（两次平行镜像＝平移，$T=2|b-a|$）是变换的合成；conccard「导数对称遗传律」（周期全遗传、对称奇偶互换）是不变量在求导下的传递；$45^\circ$ 镜面与反函数（母题二）则是「坐标互换型」的镜像。

\textbf{与通法对读}：秒判示范——$y=e^x$ 的三种变身：镜像 $e^{-x}$（值域 $(0,+\infty)$ 不变、单调反转）；旋转 $-e^{-x}$（值域 $(-\infty,0)$、单调仍增）；平移 $e^{x-h}+k$（形状不变、对称身份清零）。通法负责逐条验证，表格负责直接读取；考试书写用通法，审题预判用表格。
\logotag{高观点}

\begin{advancedonly}
\item \textbf{结构推广：任意轴／任意中心的公式、分解唯一性与变换群}。

把原点的三种变换推广到任意位置：关于直线 $x=a$ 镜像：$f(x)\to f(2a-x)$；关于点 $(a,b)$ 旋转 $180^\circ$：$f(x)\to 2b-f(2a-x)$。

\textbf{分解唯一性}：若 $f=g_1+h_1=g_2+h_2$（$g_i$ 偶、$h_i$ 奇），则 $E=g_1-g_2=h_2-h_1$ 既偶又奇，故 $E(x)=E(-x)=-E(x)$，$E\equiv0$。于是「偶部／奇部」唯一，它们是函数空间上的投影：$P_\pm^2=P_\pm$，$P_+P_-=0$——这正是「正交分解」名称的来历。

\textbf{变换群}：$M_a\circ M_b=T_{2(b-a)}$，$R_O=M_x\circ M_y$，$M^2=R^2=\mathrm{Id}$。于是「$f$ 关于 $x=a$ 与 $x=b$ 都对称 $\Rightarrow T=2|b-a|$」不过是 $M_a\circ M_b$ 作用在图象上等于恒等——双对称推周期是变换合成的推论。
\logotag{公式族}
\end{advancedonly}

\item \textbf{易错边界排雷}：诊断下面的「貌似合理」的说法，并填空防扣分铁律。

{\bfseries\color{red!70!black}错解：}因为 $f(x)+f(-x)=2x^2+2$，平均值是 $2$，所以 $f(x)=x^3+x^2+x+1$ 关于点 $(0,1)$ 中心对称；又偶部 $x^2+1$ 是偶函数，所以 $f$ 关于 $y$ 轴对称。

\ifteacher
{\bfseries\color{main}诊断：}
\begin{itemize}[itemsep=1pt]
\item \textbf{中心对称看恒等式}：$f$ 关于 $(0,b)$ 中心对称要求 $f(x)+f(-x)=2b$ 对所有 $x$ 恒成立；这里 $f(x)+f(-x)=2x^2+2$ 随 $x$ 变化，取任何 $b$ 都不是恒等式。实际上 $f$ 没有对称中心——偶部 $x^2+1$ 不是常数（偶部为常数才可能有对称中心）。
\item \textbf{部分对称≠整体对称}：$f$ 关于 $y$ 轴对称当且仅当奇部恒为 $0$；这里奇部 $x^3+x\not\equiv0$，故 $f$ 不是偶函数。分解的意义是把「可能具有不变性的部分」分离出来，而不是「每部分对称则整体对称」。
\end{itemize}

\textbf{考场防扣分铁律}：判对称只能用整体恒等式 \tk{f(-x)=\pm f(x)}（或中点公式），分解只是分离不变量的工具。
\else
{\small（先自己找出两处漏洞，再对照教师版。）}
\fi
\logotag{避坑铁律}
\end{enumerate}
\end{tcolorbox}

% ---------------------------------------------------------

\begin{tcolorbox}[
  breakable, enhanced, sharp corners,
  colback=structurecolor!3, colframe=structurecolor!70!black, boxrule=0.6pt,
  title={\bfseries 解答即笔记·五问连环母题二：$45^\circ$ 镜面与指对同构——平移型镜面 $y=x+1$ 下的「不解方程求和」},
  coltitle=black, fonttitle=\small,
  attach boxed title to top left={xshift=2mm, yshift=-2mm},
  boxed title style={sharp corners, colback=white,
    colframe=structurecolor!70!black, boxrule=0.4pt},
  top=7pt]

曲线 $C_1:y=e^x$ 关于直线 $l:y=x+1$ 的对称曲线记为 $C_2$；斜率为 $-1$ 的直线 $L:x+y=4$ 分别交 $C_1,C_2$ 于 $A(x_1,y_1)$，$B(x_2,y_2)$。求 $x_1+x_2$。

\begin{enumerate}[label=(\arabic*), itemsep=4pt]
\item \textbf{具象感知与画图}：在同一坐标系中作出 $C_1$、$C_2$ 与 $L$，标出两交点，观察弦 $AB$ 与镜面 $l$ 的关系。

\begin{center}
\begin{tikzpicture}[>=Stealth, semithick, font=\footnotesize, x=1.5cm, y=1.1cm]
  \draw[gray!22] (-1.5,-1) grid (2.6,3.5);
  \draw[->, gray] (-1.7,0)--(2.8,0) node[below right] {$x$};
  \draw[->, gray] (0,-1.2)--(0,3.7) node[above left] {$y$};
  \foreach \x in {1,2}{ \draw (\x,0.08)--(\x,-0.08) node[below=1pt] {\scriptsize \x}; }
  \foreach \y in {1,2,3}{ \draw (0.08,\y)--(-0.08,\y) node[left=1pt] {\scriptsize \y}; }
  \draw[main!70!black, very thick, domain=-1.3:1.2, samples=70, smooth] plot(\x,{exp(\x)});
  \draw[second, very thick, domain=-0.85:2.5, samples=80, smooth] plot(\x,{ln(\x+1)+1});
  \draw[third, thick, dashed, domain=-1.1:1.6, samples=40] plot(\x,{\x+1});
  \node[main!80!black] at (-1.55,1.35) {$C_1:y=e^x$};
  \node[second!80!black] at (2.55,0.55) {$C_2:y=\ln(x+1)+1$};
  \node[third!70!black] at (1.72,2.95) {$l:y=x+1$};
  \ifteacher
    \draw[red!70!black, very thick] (0.55,3.45)--(2.7,1.3);
    \node[red!70!black] at (2.75,1.18) {$L:x+y=4$};
    \fill[main] (1.0737,2.9263) circle (1.5pt) node[above left] {$A(1.074,2.926)$};
    \fill[main] (1.9263,2.0737) circle (1.5pt) node[below right] {$B(1.926,2.074)$};
    \fill[third] (1.5,2.5) circle (1.5pt) node[below left] {$M(1.5,2.5)$};
    \fill[gray] (0,1) circle (1.5pt) node[below right] {$(0,1)$};
    \node[third!70!black, align=left, font=\scriptsize] at (-1.6,3.35) {AB 中点落在 $l$ 上 $\Rightarrow$ $A,B$ 关于 $l$ 对称};
  \else
    \draw[gray!60, thick] (0.55,3.45)--(2.7,1.3);
    \node[third] at (-1.55,3.3) {描出两交点，量出中点，猜 $x_1+x_2$ 的值};
  \fi
\end{tikzpicture}
\end{center}

\ifteacher
{\small\textcolor{second}{红笔：两曲线在镜面上的公共点为 $(0,1)$（因 $e^0=1=0+1$）。数值定位 $A\approx(1.074,2.926)$，$B\approx(1.926,2.074)$，中点 $M(1.5,2.5)$ 恰在 $y=x+1$ 上——弦 $AB$ 被镜面垂直平分。猜想：$x_1+x_2=3$。}}
\else
\noteblank[3ex]
\fi
\logotag{具象感知}

\item \textbf{考纲满分通法}：四步写满，不用解超越方程。

第一步（求 $C_2$ 方程）：点关于 $y=x+b$ 的对称公式为 $(x,y)\mapsto(y-b,\,x+b)$。设 $P(x_0,y_0)\in C_2$，则 $P'(y_0-1,\,x_0+1)\in C_1$，故 $x_0+1=e^{\,y_0-1}$，即
\[
\tk{y_0=\ln(x_0+1)+1},\qquad C_2:y=\ln(x+1)+1 .
\]

第二步（列交点方程）：$A\in C_1\cap L$：\tk{e^{x_1}=y_1,\ x_1+y_1=4}；$B\in C_2\cap L$：\tk{\ln(x_2+1)+1=y_2,\ x_2+y_2=4}。

第三步（构造 $B$ 并验证）：断言 $B=(y_1-1,\,x_1+1)$。代入 $C_2$：$\ln\bigl((y_1-1)+1\bigr)+1=\ln y_1+1=x_1+1$，成立（用到 $e^{x_1}=y_1$）；代入 $L$：$(y_1-1)+(x_1+1)=x_1+y_1=4$，也成立。又 $\ln(x+1)+x$ 严格增（导数 $1+\frac1{x+1}>0$），故交点唯一，$B$ 就是它。

第四步（求和）：
\[
\tk{x_1+x_2=x_1+(y_1-1)=x_1+y_1-1=3}.
\]
\logotag{考纲通法}

\item \textbf{高观点降维秒杀}：镜面—截线垂直 $\Rightarrow$ 交点对对称 $\Rightarrow$ 坐标和公式。

定理：两条关于 $y=x+b$ 对称的曲线被直线 $x+y=m$ 所截于 $A,B$，则 $A,B$ 关于镜面对称，且
\[
\tk{x_1+x_2=m-b,\qquad y_1+y_2=m+b}
\]
（一句话推导：对称点 $(x,y)\mapsto(y-b,\,x+b)$，再代入 $A$ 满足 $x_1+y_1=m$。）代入 $b=1,m=4$：$x_1+x_2=3$，$y_1+y_2=5$。特例 $b=0$ 即退化为本章 $y=x$ 镜面的指对同构结论 $x_1+x_2=m$。

与第 (2) 问左右对读：
\begin{center}\small
\begin{tabular}{p{2.4cm}|p{5.6cm}|p{5.6cm}}
\toprule
 & 考纲通法（第 (2) 问） & 秒杀公式（第 (3) 问） \\
\midrule
思路 & 设点—验证—求和，四步，含唯一性说明 & 对称点公式 + 截线方程，两行 \\
何时用 & 解答题必须写满规范步骤 & 选填预判、答案校验 \\
风险点 & 忘记声明 $B$ 唯一；对数定义域 $x>-1$ & 忘记补偿 $b$，把 $m-b$ 写成 $m$ \\
\bottomrule
\end{tabular}
\end{center}
\logotag{高观点}

\begin{advancedonly}
\item \textbf{结构推广：另一族 $45^\circ$ 镜面与「保住函数资格」的原因}。

镜面换为 $y=-x+b$（对称点 $(x,y)\mapsto(b-y,\,b-x)$），与之垂直的截线应取斜率 $1$：$y=x+c$。同理得
\[
\tk{x_1+x_2=b-c,\qquad y_1+y_2=b+c}.
\]
特例检验：$xy=1$（关于 $y=-x$ 自对称，$b=0$）被 $y=x+c$ 所截：$x(x+c)=1\Rightarrow x^2+cx-1=0$，韦达直接给出 $x_1+x_2=-c$，与公式一致。

\textbf{为何偏偏 $k=\pm1$}：一般镜面 $y=kx+b$ 的反射会让横、纵坐标互相耦合，像曲线常跌出函数集合；只有 $k=\pm1$ 时耦合系数全部约尽，退化为干净的坐标互换（见本章 conccard「唯有 $k=\pm1$ 的镜面，能保住『函数』的资格」）。

\textbf{附注}：$C_1$ 与 $C_2$ 的唯一公共点是 $(0,1)$：令 $\varphi(x)=e^x-\ln(x+1)-1$，则 $\varphi(0)=\varphi'(0)=0$ 且 $\varphi''(x)=e^x+\frac1{(x+1)^2}>0$，故 $\varphi\ge0$、等号仅 $x=0$（即 $e^x\ge x+1$，第 6 章将用导数严格证明）。
\logotag{公式族}
\end{advancedonly}

\item \textbf{易错边界排雷}：诊断下面的「貌似合理」的解法，并填空防扣分铁律。

{\bfseries\color{red!70!black}错解：}$y=e^x$ 关于 $y=x+1$ 对称，就是把曲线向左平移 $1$ 个单位，得 $C_2:y=e^{x+1}$；再套 $y=x$ 镜面的结论，得 $x_1+x_2=4$。

\ifteacher
{\bfseries\color{main}诊断：}
\begin{itemize}[itemsep=1pt]
\item \textbf{对称≠平移}：关于 $y=x+1$ 的对称是「坐标角色互换」$(x,y)\mapsto(y-1,\,x+1)$，不是同向加减。镜面上的点 $(0,1)$ 必在 $C_2$ 上，而错写的 $y=e^{x+1}$ 在 $x=0$ 处给 $e\ne1$，立即矛盾。
\item \textbf{公式漏补偿}：即便曲线求对，$y=x$ 的结论也不能直接搬到 $y=x+1$，须用 $x_1+x_2=m-b=3$。数值复核：$x_1\approx1.074$，$x_2\approx1.926$，和为 $3$ 而非 $4$。
\end{itemize}

\textbf{考场防扣分铁律}：先用替换口径 \tk{F(y-b,\,x+b)=0} 把对称曲线方程写准，再谈任何结论；求和公式里的 $b$ 不能丢。
\else
{\small（先自己找出两处漏洞，再对照教师版。）}
\fi
\logotag{避坑铁律}
\end{enumerate}
\end{tcolorbox}

% ---------------------------------------------------------

\begin{notezone}{章末归档：对称即不变性 + $45^\circ$ 镜面求和公式族}
\textbf{对称即不变性卡}：镜像 $f(x)\to f(-x)$ 保持值域与奇偶身份、反转单调方向；旋转 $f(x)\to-f(-x)$ 保持单调方向与奇偶身份、翻转值域；平移保持形状、破坏对称位置。任意 $f=\dfrac{f(x)+f(-x)}{2}+\dfrac{f(x)-f(-x)}{2}$：偶部是镜像不变量、奇部是旋转不变量，分解唯一（$P_\pm^2=P_\pm$，$P_+P_-=0$）。

\textbf{$45^\circ$ 镜面求和公式族}：$y=x+b$ 镜面 × 截线 $x+y=m$ $\Rightarrow$ $x_1+x_2=m-b$，$y_1+y_2=m+b$；$b=0$ 即指对同构经典款；$y=-x+b$ 镜面 × 截线 $y=x+c$ $\Rightarrow$ $x_1+x_2=b-c$，$y_1+y_2=b+c$。

\textbf{五问与笔记}：(1) 具象感知与画图 $\to$ 概念引入笔记；(2) 考纲满分通法 $\to$ 规范书写笔记；(3) 高观点降维秒杀 $\to$ 大招与二级结论笔记；(4) 结构推广与变式 $\to$ 公式族笔记；(5) 易错边界排雷 $\to$ 错题本笔记。
\end{notezone}
